\section{Implementation}\label{sec:impl}

The implementation is one repository; Perch is a pinned dependency used
unchanged, and the circuit is compiled and proven with exactly Perch's pinned
toolchain, nargo 1.0.0-beta.9 and Barretenberg 0.87.0, so its proofs verify
with Perch's Soroban verifier code\claim{circuit-size}. The optimized
WebAssembly is 33{,}479 bytes for the guardian, 68{,}269 for the registry,
45{,}616 for the verifier and 10{,}152 for the enclave-image
registry\claim{contract-sizes}.

\paragraph{Circuit.} The circuit is a library of components (RSA-2048
PKCS\#1 v1.5 with SHA-256; the structural scan of the signed block; the
\texttt{From}, \texttt{DKIM-Signature} and \texttt{Subject} parsers; packing;
the registry leaf and Merkle root) composed by one relation function. It
vendors only noir-bignum, for the 2048-bit modular multiplication, and
Noir's SHA-256 library, each pinned file by file; Poseidon2 is Perch's
vendored copy; no parser is reused\claim{circuit-vendoring}. The header parsers are written against
the grammars of \S\ref{sec:design-circuit}, and every position in the header
is an unconstrained hint that constrained code re-checks. At a 1024-byte
header bound the same relation has 189{,}503 gates and the same padded size;
the 1536-byte bound costs the prover more but leaves the verification key
and the on-chain cost unchanged\claim{circuit-size}. That bound admits every
launch domain's recovery email and, with a 124-byte address, a display name
of 536 bytes at the tightest provider\claim{header-bound-fits}.

\paragraph{Contracts.} The guardian, the registry, the verifier and the
enclave-image registry are Soroban contracts in Rust. The verifier is a
constructorless, immutable contract that embeds Perch's UltraHonk verifier
from the same Perch commit as the other Perch dependencies, with the email
circuit's verification key compiled in, and reports the key's SHA-256 as its
circuit id; the guardian refuses at construction a verifier whose circuit id
differs from its own configuration\claim{guardian-construction}. A real
proof of each launch domain's recovery email verifies on chain through it,
is refused with any public input changed, and counts as exactly one guardian
vote\claim{onchain-email-vote}. The
registry stores its tree with the host's Poseidon2 sponge and Perch's table
of empty-subtree hashes, so the root it stores is the root the circuit proves
against\claim{registry-hash-parity}.

\paragraph{Wallet.} The browser side is TypeScript: an encoder from a raw
message to the witness, the prover worker (noir\_js 1.0.0-beta.9 and bb.js
0.87.0), which checks its downloads against pinned digests and verifies each
proof locally against the guardian's circuit id\claim{browser-prover}, the
enrollment and recovery pages, and the key lookup.

\paragraph{Alert service.} The host watches the controller over Stellar
RPC with a stored cursor, retries failed alerts, and never repeats one; the
vault runs as a separate process standing in for the enclave. The wallet's
HPKE reproduces RFC~9180's test vector, and the vault's Web Push encryption
RFC~8291's example\claim{alert-formats}.

\paragraph{Mocks.} Every third party is a test-only mock in its published
format: the mail providers sign with test-only keys in each provider's layout,
established from public mailing-list archives\claim{provider-mock-formats};
DNS answers follow RFC~6376's key records and the DNS-over-HTTPS JSON
format\claim{key-lookup}; the sending services follow their published APIs,
Web Push its RFCs, and the enclave's attestation AWS's documented
format~\cite{aws-nitro-attestation}\claim{vault-custody}; the DNS attester
signs test-only Ed25519 attestations.
